\section{Organización del documento}
\label{sec:intro_organizacion}
Después de esta introducción, el documento se organiza en seis capítulos:

\begin{description}
  \item[Capítulo~\ref{cp:marco}, Marco teórico.] Este capítulo presenta los conceptos necesarios para leer el
  resto del documento y revisa los trabajos más cercanos.
  \item[Capítulo~\ref{cp:metodos}, Materiales y métodos.] Este capítulo presenta los conjuntos de datos, los
  modelos, el \emph{pipeline} que encadena la clasificación y la generación del reporte, el protocolo de
  evaluación y el diseño de los experimentos.
  \item[Capítulo~\ref{cp:res_clasif}, Resultados de clasificación.] Este capítulo presenta la comparación de los
  clasificadores supervisados y de los \glspl{vlm}, la detección de melanoma y la evaluación del clasificador elegido
  con imágenes de otro conjunto de datos y con imágenes degradadas, junto con su calibración y su costo, y la comparación con el estado del arte.
  \item[Capítulo~\ref{cp:res_gen}, Resultados de generación.] Este capítulo presenta la evaluación de los reportes
  a lo largo del \emph{pipeline}: el generador sin diagnóstico, la cascada con el clasificador, el ruteo por
  concordancia, la comparación entre generadores, sus costos computacionales y el contraste con el estado del arte.
  \item[Capítulo~\ref{cp:discusion}, Discusión.] Este capítulo presenta la interpretación conjunta de los
  resultados y las limitaciones del trabajo.
  \item[Capítulo~\ref{cp:conclusiones}, Conclusiones.] Este capítulo presenta la respuesta a las preguntas
  planteadas, los hallazgos principales, el
  balance de los objetivos, los aportes del trabajo y el trabajo futuro.
\end{description}

Por último, dos apéndices reúnen el material que, por ser muy detallado, no se incluye en el cuerpo de los
capítulos: los \emph{prompts} y las plantillas que usan los modelos, y la información necesaria para reproducir los
experimentos.
